\section{Guarantees and Correctness Properties}

The Bailout Protocol is not merely a design pattern; it is a formally specified subsystem with three independently verifiable correctness properties. Each property admits a structural proof sketch grounded in the BX3 architecture rather than empirical observation. We establish these properties in sequence: SAFETY, LIVENESS, and ACCOUNTABILITY.

\subsection{SAFETY: No Autonomous Collapse}

\textbf{Theorem (Safety).} \emph{Under the Bailout Protocol, no node can transition to a state in which all resolution paths lead exclusively through Machine Accountability Anchors.}

\begin{proof}[Proof Sketch]
The BX3 architecture enforces a strict containment property: the Purpose layer of every node is anchored to a Human Accountability Anchor (HAA) at creation. This anchoring is immutable — the Root Pipe Protocol permits authority projection downward, but no mechanism permits upward escape from human anchoring.

Consider an arbitrary node $N$ at some timestep $t$. Let $S_t$ denote the set of active resolution paths available to $N$. By construction, $S_t$ always contains at least one path $p_{\text{human}}$ that terminates at an HAA, because:
\begin{enumerate}
    \item The node's Worksheet embeds the complete Bailout Protocol at initialization.
    \item The three trigger conditions (Capability Boundary, Safety Envelope Violation, Accountability Boundary) are checked before any Machine Actor closure.
    \item Machine Actor exclusion is structurally enforced: propagation logic explicitly filters out Machine Accountability Anchors from the candidate resolution set.
\end{enumerate}

An \emph{autonomous collapse} event would require $p_{\text{human}} = \emptyset$. This would imply either: (a) the HAA anchoring was removed (contradicted by immutability of the Purpose layer), or (b) the trigger conditions were bypassed (contradicted by the hardcoded nature of trigger evaluation in the Fact Layer). Therefore $p_{\text{human}} \neq \emptyset$ for all $t$, and SAFETY holds. \hfill $\square$
\end{proof}

\textbf{Consequences.} SAFETY implies that human oversight is not a optional feature or a secondary fallback — it is a structural invariant. Even in the degenerate case where every Machine actor in a resolution chain rejects the exception, the propagation terminates at an HAA. The system cannot route an unresolved exception to \texttt{null} or discard it silently.

\subsection{LIVENESS: Eventual Human Contact}

\textbf{Theorem (Liveness).} \emph{Any exception event generated by a node under the Bailout Protocol reaches a Human Accountability Anchor within a bounded number of propagation steps.}

\begin{proof}[Proof Sketch]
Define the \emph{depth} of a node as its graph-theoretic distance from the Human Root in the BX3 hierarchy. Let $d(N)$ denote the depth of node $N$, where $d(\text{Human Root}) = 0$. The maximum depth of any node in the system is $D_{\max}$, a finite constant determined by the provisioned hierarchy.

The Bailout Protocol trigger conditions satisfy the following:
\begin{enumerate}
    \item \textbf{Finite branching:} Each node propagates exceptions to at most one parent per event.
    \item \textbf{Upward monotonicity:} Propagation always moves to a parent of strictly lower depth ($d_{\text{parent}} < d_{\text{child}}$).
    \item \textbf{No cycles:} The hierarchy is a rooted tree by construction; cycles are architecturally impossible.
\end{enumerate}

Therefore, any exception originating at depth $d$ propagates at most $d$ steps before reaching depth $0$ (Human Root). The worst-case latency bound is $O(D_{\max})$ propagation steps. Since $D_{\max}$ is finite and constant for a given deployment, human contact is guaranteed within a bounded number of steps, satisfying LIVENESS. \hfill $\square$
\end{proof}

\textbf{Consequences.} LIVENESS does not guarantee \emph{speed} — only \emph{eventual contact}. The bound $D_{\max}$ depends on hierarchy depth; a deeply nested deployment may require dozens of propagation steps before an HAA is reached. However, no exception can circulate indefinitely or be absorbed by the Machine actor layer. The protocol is starvation-free with respect to human review.

\subsection{ACCOUNTABILITY: Human Always Identifiable}

\textbf{Theorem (Accountability).} \emph{For any event processed by the Bailout Protocol, there exists a unique, log-identifiable Human Accountability Anchor who bears formal responsibility for the outcome.}

\begin{proof}[Proof Sketch]
Define the \emph{Accountability Chain} as the ordered sequence of nodes that process a given exception from origin to resolution. Let $C = (e_0, a_1, a_2, \dots, a_n, h)$ denote this chain, where $e_0$ is the originating node, each $a_i$ is a Machine Accountability Anchor (if any), and $h$ is the terminal HAA.

The Forensic Ledger records three immutable facts for every event:
\begin{enumerate}
    \item The identity of the HAA $h$ that issues the terminal directive.
    \item The trigger condition classification at each propagation step.
    \item The full context package forwarded at each step, including Confidence Score.
\end{enumerate}

The Ledger is cryptographically sealed across three planes (Purpose, Bounds Engine, Fact Layer) using SHA-256 forensic sealing. This ensures that retroactively identifying the responsible HAA requires only a read of the sealed event record — no inference, no reconstruction, no ambiguity.

Critically, the Machine Actor Exclusion property ensures that no Machine Accountability Anchor can appear in the terminal position of $C$. The chain always terminates at a human. The unique HAA identity is therefore guaranteed by construction. \hfill $\square$
\end{proof}

\textbf{Consequences.} ACCOUNTABILITY ensures that no autonomous action can be legally or operationally orphaned. Even in cases of network partition, node failure, or deliberate sabotage, the Forensic Ledger preserves an unbroken chain connecting every event to a human who bears responsibility. This is the structural answer to the ``black box'' problem that conventional autonomous systems cannot resolve.

\subsection{Interdependency of the Three Properties}

The three properties are not independent axioms — they form a logically coherent constraint system:

\begin{itemize}
    \item \textbf{SAFETY} $\Rightarrow$ the human anchor exists and is reachable.
    \item \textbf{LIVENESS} $\Rightarrow$ the anchor is reached within bounded steps.
    \item \textbf{ACCOUNTABILITY} $\Rightarrow$ the anchor is uniquely and verifiably identified.
\end{itemize}

Together, these establish the \emph{Bailout Invariant}: \emph{for any exception event in a BX3 system, there exists exactly one identifiable human who bears accountability, and this human is reachable within a bounded, finite number of propagation steps.}

This invariant holds regardless of the state of the Machine actor layer, the depth of the recursive hierarchy, or the presence of network partitions. It is a structural guarantee, not a statistical one.

\subsection{Discussion: Relation to Classical Safety-Liveness Framework}

Classical distributed systems theory distinguishes between \emph{safety properties} (``nothing bad ever happens'') and \emph{liveness properties} (``something good eventually happens''). The Bailout Protocol's three properties map cleanly onto this framework:

\begin{itemize}
    \item \textbf{SAFETY} is a safety property: autonomous collapse is a permanent, irreversible bad outcome, and the theorem guarantees it never occurs.
    \item \textbf{LIVENESS} is a liveness property: human contact is a good outcome that must eventually occur, and the bounded propagation guarantee provides this.
    \item \textbf{ACCOUNTABILITY} is neither purely safety nor purely liveness — it is a \emph{knowledge property} ensuring that responsibility is not just achieved but \emph{attributed}, closing the loop between action and identification.
\end{itemize}

The addition of ACCOUNTABILITY as a third axis distinguishes the Bailout Protocol from conventional exception-handling frameworks, which typically address only safety and liveness without mechanisms for unique human identification in recursive autonomous systems.

\end{document}